\begin{table}[htbp]
\centering
\caption{Summary Statistics: Pre-Reform Period (2012--2018)}
\label{tab:summary}
\begin{tabular}{lcc}
\hline\hline
 & Mining Districts & Non-Mining Districts \\
\hline
\addlinespace
\multicolumn{3}{l}{\textit{Panel A: District Characteristics}} \\
\addlinespace
Number of districts & 11 & 104 \\
\addlinespace
\multicolumn{3}{l}{\textit{Panel B: Nighttime Light Intensity (Pre-Reform)}} \\
\addlinespace
Mean NTL (nW/cm$^2$/sr) & 0.566 & 0.267 \\
SD NTL & 0.681 & 2.304 \\
Mean asinh(NTL) & 0.465 & 0.075 \\
SD asinh(NTL) & 0.490 & 0.393 \\
\addlinespace
\multicolumn{3}{l}{\textit{Panel C: Nighttime Light Intensity (Post-Reform)}} \\
\addlinespace
Mean NTL (nW/cm$^2$/sr) & 0.558 & 0.246 \\
SD NTL & 0.648 & 2.103 \\
\addlinespace
Observations (district-years) & \multicolumn{2}{c}{1380} \\
Years & \multicolumn{2}{c}{2012--2023} \\
\hline\hline
\end{tabular}
\begin{tablenotes}
\item \textit{Notes:} Mining districts are those in Zambia's Copperbelt and North-Western provinces with active copper/cobalt mining operations: Kitwe, Ndola, Mufulira, Chingola, Luanshya, Kalulushi, Chililabombwe, Mpongwe, Solwezi, Kalumbila, and Kasempa. NTL is the annual mean nighttime light intensity from NASA VIIRS Black Marble (VNP46A4) measured in nanowatts per square centimeter per steradian.
\end{tablenotes}
\end{table}

